\section{A polynomial-size quantifier-free decision backend}\label{sec:qf}
The general Presburger construction supports arbitrary finite-fibre source
predicates and complete target equality. For the modal schemas in the
finite-defect cover, a narrower syntax admits a polynomial-size existential
encoding. This provides an executable decision backend and a second route
to the small-witness conclusion.

Quantifier-free encodings of counting properties on flat runs already appear
in Decker--Pirogov~\cite{flat}. Here the additional obligation is to encode
all spatial coincidences and chronologically earlier flips in a complete
mutable field. The restricted strides below make those parity constraints
polynomial in the supplied schema size.

Let
\[
 \mathcal D=\{(\pm2,\pm2),(\pm2,0),(0,\pm2),(\pm4,0),(0,\pm4)\}.
\]
Every powered word must restore heading and have full drift in $\mathcal D$.
The source is a finite XOR of segments $a+jv$, $0\le j<L$, with
$v\in\mathcal D$, and singletons. Anchors, lengths and exponents may be
affine expressions in shared integer parameters; lengths and exponents
are constrained nonnegative. Explicit $m$-point inputs additionally require
pairwise distinct coordinates. No coordinate or repetition bound is imposed.

\begin{theorem}[Quantifier-free physical realization]
For this syntax, a polynomial-time construction produces a quantifier-free
linear-integer/Boolean circuit whose existential satisfiability is equivalent
to literal execution of the complete prescribed word. Exact contained-defect
counts, the final defect flag, endpoint time and pose, and any polynomial
number of endpoint-colour queries have polynomial-size circuits as well.
\end{theorem}
\begin{proof}
Represent each phase of a power by the injective read stream
$a_i+j d$, $0\le j<k$, with its fixed required colour. Its actual colour
is the XOR of the source and every earlier planned write at that point.
For streams in different blocks, temporal order is fixed. Within one
power of period $P$, a phase-$r$ occurrence at index $j+s$ precedes
phase $i$ at index $j$ exactly when $Ps+r<i$. This retains the intra-loop
chronology rather than grouping all occurrences by phase.

For a source or write stream $b+\ell v$, solve
$a+jd=b+\ell v$ with its finite index bounds. Nonparallel strides give at
most one hit by the determinant formula. Parallel strides give a guarded
arithmetic interval in $j$ with modulus
$|v_\alpha|/\gcd(|d_\alpha|,|v_\alpha|)$, which is always 1 or 2 for
$\mathcal D$. Splitting $j=2z+e$ leaves a guarded ordinary interval.
Floors and exact divisibility use only fixed integer divisors and therefore
have constant-size linear defining constraints.

On a target residue interval $[0,K]$, the XOR of these interval contributions
is a step function. It has the required constant colour exactly when its
value at zero is correct and every active interior boundary has even total
jump parity. Check this at each candidate interval start and one-past-end,
including all coincident boundaries. For $q$ contributions, the equality
and parity tests use $O(q^2)$ shared wire operations; no ordering of symbolic
endpoints is enumerated. With $R$ read streams and $M$ source terms the
parity stage has size $O(R(R+M)^2)$, apart from shared affine expressions.
Induction through the finite planned word proves that these colour tests
are sufficient as well as necessary for physical execution.

For lexical counts, retain the preceding at most nine turns, their length,
the defect count and the final-event flag as shared registers. For a power
$u^k$, at most $\lceil9/|u|\rceil$ copies precede suffix stabilization;
later copies contribute a fixed cyclic defect count. Constant-size guarded
transfers select the short cases and the stabilized case. The $k=0$ branch
preserves the entire incoming state. Sharing these transfers between blocks
gives a circuit linear in spelled length and block count, instead of a
Cartesian list of zero/positive choices.

Finally, an endpoint-colour query XORs initial membership with membership
in every read stream. This retains all unqueried remote cells. Named gates
and linear definitions preserve polynomial size throughout and prove the
theorem.
\end{proof}

\begin{corollary}[Unary budgets with binary side conditions]
For explicit $m$ coordinate pairs, unary $N\ge1$ and a quantifier-free
linear-integer side condition $Q$ of binary length $b$, existence of a
complete $m$-black source whose $N$th contained defect satisfies $Q$ is
in NP. Positive witnesses have polynomial bit length in $m+N+b$.
For unary $N$, the minimum initial cost $\mu(N)$ is in $\mathrm{FP}^{\mathrm{NP}}$.
\end{corollary}
\begin{proof}
The inherited complete cover supplies a schema with at most $5(N+1)$
powers and $2048(N+1)+N$ fixed letters. Its powered bodies are phases of
$W,T,A$, with strides in $\mathcal D$. Guess this polynomial-length
schema, apply the theorem to the $m$ distinct source points, and require
exactly $N$ defects with a bad last letter and condition $Q$.
Guessing Boolean branches reduces the circuit to polynomially many linear
integer constraints. The integer small-solution bound used in
Section~\ref{sec:small} supplies polynomial-bit witnesses, which can be
checked in polynomial time. This uses the complete physical cover;
satisfiability of a graph relaxation alone would not suffice.

Theorem~\ref{thm:six} gives $\mu(N)\le6$. A finite prefix can be padded
by unread black cells outside its footprint, so the NP query for exactly $m$
points decides $\mu(N)\le m$ without additional side conditions. Searching
$m\in\{0,\ldots,6\}$ uses a constant number of NP-oracle queries, each of
polynomial unary-$N$ size, and proves the stated function-class bound.
\end{proof}

The encoding applies to the stated segment-source syntax and endpoint
colour queries. The general compiler handles quantified source predicates,
zero-displacement powers and complete-field target equality. Polynomial circuit size does not
imply a polynomial deterministic solver, and the NP statement is not a
claim for binary-encoded $N$.

The reproduction archive checks the encoding by exact intersection
identities and complete finite endpoint comparisons. Its Z3 and ISL/GMP
backends are part of the computational evidence described in the
reproduction section.
